% 《大学物理简明教程》第5版，第4章解答题4.3.11—4.3.17。
% 来源印刷/PDF第104—105页，波形图模型按原图定位；待学术二校。
\subsection{解答题：平面行波与相位}

\sourceproblem{4.3.11}{104}{104}
\problemtext 已知沿$x$轴传播的平面简谐波
$y=A\cos(Bt-Cx)$，$A,B,C>0$。求振幅、波速、频率、周期和波长；
原点的振动方程；以及沿传播方向相距$\lambda/2$两点的相位差。
\workings
与标准式$y=A\cos(\omega t-kx)$比较，$\omega=B,k=C$，故
\[
\boxed{A=A,\quad u=\frac BC,\quad
\nu=\frac B{2\pi},\quad T=\frac{2\pi}{B},\quad
\lambda=\frac{2\pi}{C}}.
\]
$x=0$处振动为$\boxed{y_O(t)=A\cos Bt}$。
同一时刻，相距$\lambda/2$的相位差大小
\[
\boxed{|\Delta\varphi|=C\lambda/2=\pi}.
\]
\crosscheck 固定相位$Bt-Cx=\mathrm{const}$的轨迹$x=(B/C)t+\mathrm{const}$，方向为$x$正向。

\sourceproblem{4.3.12}{104}{104}
\problemtext 沿细绳传播的简谐波
$y=0.05\cos(10\pi t-4\pi x)$（SI）。求：
(1) 各质点最大速度和加速度；
(2) $x=0.2\,\mathrm m$处$t=1\,\mathrm s$时的相位与原点在哪个时刻相同；
(3) 该相位在$t=1.25\,\mathrm s$时传播到哪里。
\workings
$\omega=10\pi\,\mathrm{rad/s},k=4\pi\,\mathrm{rad/m}$，
$u=\omega/k=2.5\,\mathrm{m/s}$。
\[
\boxed{v_{\max}=A\omega=0.5\pi\,\mathrm{m/s}},\qquad
\boxed{a_{\max}=A\omega^2=5\pi^2\,\mathrm{m/s^2}}.
\]
$x=0.2,t=1$的连续相位$\theta=10\pi-0.8\pi=9.2\pi$ rad
（模$2\pi$为$1.2\pi$）。原点有相同相位时
$10\pi t_0=9.2\pi+2n\pi$，故
\[
\boxed{t_0=0.92+0.2n\ \mathrm s,\quad n\in\mathbb Z}.
\]
追踪该同一相位的波峰/相位面到$t=1.25$，
\[
\boxed{x=0.2+u(1.25-1)=0.825\,\mathrm m}.
\]
\crosscheck 传播距离为$2.5\times0.25=0.625$ m。相位相同的\emph{其他周期副本}的坐标相隔整波长$\lambda=0.5$ m。

\sourceproblem{4.3.13}{104--105}{104--105}
\problemtext 原题4.3.13的某时刻空间波形沿$x$轴：$O$处过零向上，$A$为第一个正峰，$B$为随后的向下过零点，$C$为再下一个向上过零点。分别讨论波沿$+x$和$-x$传播时该时刻四点各自振动的相位。
\workings
以$A$到$C$间图像确定各点的\emph{空间}相位：
$\theta_s(O)=0,\theta_s(A)=\pi/2,\theta_s(B)=\pi,\theta_s(C)=2\pi$，故空间函数可写$y(x,0)=A_0\sin kx$（起点作适当选择）。
若波沿$+x$传播，
$y=A_0\cos(\omega t-kx+\pi/2)$；
在该时刻，以余弦式定义的四点时间相位为
\[
\boxed{\varphi_+(O)=\frac\pi2,\quad
\varphi_+(A)=0,\quad
\varphi_+(B)=-\frac\pi2,\quad
\varphi_+(C)=-\frac{3\pi}{2}}.
\]
若波沿$-x$传播，
$y=A_0\cos(\omega t+kx-\pi/2)$，
相应
\[
\boxed{\varphi_-(O)=-\frac\pi2,\quad
\varphi_-(A)=0,\quad
\varphi_-(B)=\frac\pi2,\quad
\varphi_-(C)=\frac{3\pi}{2}}.
\]
\crosscheck 相位仅在模$2\pi$意义下确定；本解固定采用\emph{时间余弦相位}，避免把波形的空间正弦相位直接当成质点的时间相位。

\sourceproblem{4.3.14}{104--105}{104--105}
\problemtext 沿$x$轴正方向传播的平面简谐波，波速$u=5\,\mathrm{m/s}$，波长$\lambda=2\,\mathrm m$。原点的振动曲线在$t=0$过零并向正位移运动，振幅$0.1\,\mathrm m$。求波函数、$t=0$的波形及距原点$0.5\,\mathrm m$处的振动曲线。
\workings
$f=u/\lambda=2.5\,\mathrm{Hz}$，
$\omega=5\pi\,\mathrm{rad/s}$，$k=\pi\,\mathrm{rad/m}$。
原点初相以正弦表述为零且初速为正，
\[
\boxed{y(x,t)=0.1\sin(5\pi t-\pi x)\,\mathrm m}.
\]
$t=0$时$\boxed{y(x,0)=-0.1\sin(\pi x)\,\mathrm m}$；
$x=+0.5$ m点的振动方程
\[
\boxed{y(0.5,t)=0.1\sin(5\pi t-\pi/2)
=-0.1\cos(5\pi t)\,\mathrm m}.
\]
\crosscheck 右行波在$t=0$空间原点的正斜率是负的，和原点随时间正速度相反：
$\partial_ty=-u\,\partial_xy$。

\sourceproblem{4.3.15}{105}{105}
\problemtext 图4.3.15给出$t=0$与$t=0.5\,\mathrm s$的波形(a)(b)，波沿$x$正方向传播。振幅$0.1\,\mathrm m$，(a)中相邻正峰间距约$4\,\mathrm m$，(b)相对(a)右移$1\,\mathrm m$。求波函数和$x=1\,\mathrm m$的$P$点振动方程。
\workings
由图可得$\lambda=4\,\mathrm m$、$k=\pi/2\,\mathrm{rad/m}$，
移速$u=1/0.5=2\,\mathrm{m/s}$，$\omega=ku=\pi\,\mathrm{rad/s}$。
$t=0$曲线在$x=0$向上穿过零，故
\[
\boxed{y(x,t)=0.1\sin\left(\frac\pi2x-\pi t\right)\,\mathrm m}.
\]
代入$x=1$，
\[
\boxed{y_P(t)=0.1\cos(\pi t)\,\mathrm m}.
\]
\crosscheck 源图只画有限波段，位移$1$ m取为向右的\emph{最小正}图示平移；相差整波长的平移给出等效波形，但由此产生的传播速度不唯一，须以题图暗示的相邻状态识别。

\sourceproblem{4.3.16}{105}{105}
\problemtext 波沿$x$正方向以$u=10\,\mathrm{m/s}$传播，$\lambda=4\,\mathrm m$，$t=0$波形如图4.3.16，振幅$A=0.1\,\mathrm m$，原点$y(0,0)=0.05$ m且曲线向上，$P$点处在负半振幅$y_P=-0.05$ m的\emph{上升空间分支}。求波函数、$P$点振动方程、其坐标及回到平衡位置的最短时间。
\workings
$k=2\pi/\lambda=\pi/2\,\mathrm{rad/m}$，
$\omega=uk=5\pi\,\mathrm{rad/s}$。
令初始空间函数$y(x,0)=A\sin(kx+\varphi)$，
由$y(0,0)=A/2$且空间斜率为正，得$\varphi=\pi/6$。
\[
\boxed{y(x,t)=0.1\sin\left(\frac\pi2x-5\pi t+\frac\pi6\right)\,\mathrm m}.
\]
在$P$点，空间相位为$11\pi/6$（因$\sin=-1/2$且空间斜率为正），选离原点图示的第一个此类点：
\[
\frac\pi2x_P+\frac\pi6=\frac{11\pi}{6},
\quad
\boxed{x_P=\frac{10}{3}\,\mathrm m\approx3.333\,\mathrm m}.
\]
故
\[
\boxed{y_P(t)=0.1\sin\left(\frac{11\pi}{6}-5\pi t\right)\,\mathrm m}.
\]
$t=0$时该质点速度为负，需先经过负极值再到零。时间相位从$11\pi/6$下降至$\pi$时首次过零，变化$5\pi/6$，
\[
\boxed{t_{\min}=\frac{5\pi/6}{5\pi}=\frac16\,\mathrm s}.
\]
\crosscheck 该图形题给出的是\emph{选定的第一个上升空间分支}；若图中的$P$并非这一分支，坐标与最短时间须按实际斜率重新判断，故仍标记\textbf{待图示二校}。

\sourceproblem{4.3.17}{105}{105}
\problemtext 给定平面简谐波中$P$点振动方程
$y_P=A\cos(\omega t+\varphi_0)$；分别在原图(a)中以$O$为原点、$P$位于$x=l$且波沿$+x$传播，和图(b)中以$P$为原点、$Q$位于$x=b$且波沿$-x$传播时求波函数及$Q$点振动方程。
\workings
(a) 右行波按参考点$P$延迟时间$(x-l)/u$：
\[
\boxed{y_a(x,t)=A\cos\left[\omega t+\varphi_0-\frac{\omega}{u}(x-l)\right]}.
\]
若$Q$在$P$右边$b$处，则
$\boxed{y_{Q,a}(t)=A\cos(\omega t+\varphi_0-\omega b/u)}$。
(b) 以$P$为$x=0$，左行波使用$t+x/u$：
\[
\boxed{y_b(x,t)=A\cos\left(\omega t+\varphi_0+\frac{\omega x}{u}\right)},
\]
\[
\boxed{y_{Q,b}(t)=A\cos(\omega t+\varphi_0+\omega b/u)}.
\]
\crosscheck 对两式分别求等相位轨迹，(a)有$dx/dt=+u$、(b)有$dx/dt=-u$；代入$P$位置均恢复给定振动方程。
